% !TEX program = xelatex
\documentclass[10pt,aspectratio=169]{beamer}

\usepackage[T1]{fontenc}
\usepackage{lmodern}

% Chinese (Simplified) support for the translated deck. xeCJK routes CJK
% codepoints to a CJK font while leaving Latin text/math on the fontspec
% faces set per-deck; this is lighter-touch than full `ctex`, which fights
% beamer/metropolis's own fontspec setup. Requires XeLaTeX (already the
% engine for these decks) and Noto CJK SC (installed system-wide).
\usepackage{xeCJK}
\setCJKmainfont{Noto Sans CJK SC}
\setCJKsansfont{Noto Sans CJK SC}
\setCJKmonofont{Noto Sans Mono CJK SC}

% Header bar matches the rafael-sharon toronto-may-2026 reference deck:
% miniframes navigation with subsection ticks, dark-teal section/subsection
% bars at the top of every content frame.
\definecolor{bg_subsec}{RGB}{43, 66, 70}
\definecolor{bg_sec}{RGB}{51, 77, 83}

\usetheme{metropolis}
\useoutertheme[subsection=true]{miniframes}
\setbeamercolor{subsection in head/foot}{fg=white, bg=bg_subsec}
\setbeamercolor{section in head/foot}{fg=white, bg=bg_sec}

\usepackage{appendixnumberbeamer}

\usepackage{amsmath, amssymb, amsthm}
\usepackage{mathtools}
\usepackage{bm}
\usepackage{graphicx}
\usepackage{booktabs}

% Code listings
\usepackage{listings}
\usepackage{xcolor}
\definecolor{jl-keyword}{HTML}{8959A8}
\definecolor{jl-string}{HTML}{718C00}
\definecolor{jl-comment}{HTML}{8E908C}
\lstdefinelanguage{Julia}{
  keywords={if,else,elseif,end,for,while,function,return,using,import,
            module,struct,abstract,type,export,begin,let,do,in,true,false,
            const,where},
  morecomment=[l]{\#},
  morestring=[b]",
  sensitive=true,
}
\lstset{
  basicstyle=\ttfamily\footnotesize,
  language=Julia,
  keywordstyle=\color{jl-keyword}\bfseries,
  stringstyle=\color{jl-string},
  commentstyle=\color{jl-comment},
  showstringspaces=false,
  breaklines=true,
  tabsize=2,
  captionpos=b,
  frame=single,
  framerule=0pt,
  backgroundcolor=\color{black!3},
  extendedchars=true,
  literate=
    {β}{{$\beta$}}1 {α}{{$\alpha$}}1 {γ}{{$\gamma$}}1 {δ}{{$\delta$}}1
    {ε}{{$\varepsilon$}}1 {ζ}{{$\zeta$}}1 {η}{{$\eta$}}1 {θ}{{$\theta$}}1
    {κ}{{$\kappa$}}1 {λ}{{$\lambda$}}1 {μ}{{$\mu$}}1 {ν}{{$\nu$}}1
    {ξ}{{$\xi$}}1 {π}{{$\pi$}}1 {ρ}{{$\rho$}}1 {σ}{{$\sigma$}}1
    {τ}{{$\tau$}}1 {φ}{{$\varphi$}}1 {χ}{{$\chi$}}1 {ψ}{{$\psi$}}1
    {ω}{{$\omega$}}1 {Λ}{{$\Lambda$}}1 {Δ}{{$\Delta$}}1 {Σ}{{$\Sigma$}}1,
}

% Common math macros (kept in sync with paper/main.tex)
\newcommand{\R}{\mathbb{R}}
\newcommand{\E}{\mathbb{E}}
\newcommand{\Prob}{\mathbb{P}}
\newcommand{\norm}[1]{\left\| #1 \right\|}
\newcommand{\Vstar}{V^{\star}}

\newcommand{\p}[1]{\left( #1 \right)}
\newcommand{\psq}[1]{\left[ #1 \right]}

% Section page: use the long-form section name (\insertsection) instead of
% the short-form (\insertsectionhead) so \section[short]{long} renders
% `long` on the section-title slide and `short` in the header. Cloned
% from metropolis's `progressbar` section page template.
\setbeamertemplate{section page}{%
  \centering%
  \begin{minipage}{22em}%
    \raggedright%
    \usebeamercolor[fg]{section title}%
    \usebeamerfont{section title}%
    \insertsection\\[-1ex]%
    \usebeamertemplate*{progress bar in section page}%
    \par%
    \ifx\insertsubsection\@empty\else%
      \usebeamercolor[fg]{subsection title}%
      \usebeamerfont{subsection title}%
      \insertsubsection%
    \fi%
  \end{minipage}%
  \par%
  \vspace{\baselineskip}%
}

% Project-specific macros
\newcommand{\Vstart}{V^{\mathrm{start}}}
\newcommand{\Vend}{V^{\mathrm{end}}}
\newcommand{\Vstarttil}{\widetilde{V}^{\mathrm{start}}}
\newcommand{\Vendtil}{\widetilde{V}^{\mathrm{end}}}
\newcommand{\Vstarthat}{\widehat{V}^{\mathrm{start}}}
\newcommand{\Vendhat}{\widehat{V}^{\mathrm{end}}}
\newcommand{\Lstart}{\Lambda^{\mathrm{start}}}
\newcommand{\Lend}{\Lambda^{\mathrm{end}}}
\newcommand{\xstart}{x^{\mathrm{start}}}
\newcommand{\xend}{x^{\mathrm{end}}}
\newcommand{\CV}{\mathcal{V}}
\newcommand{\CL}{\mathcal{L}}


\usepackage{fontspec}

\setmainfont{Latin Modern Roman}
\setsansfont{Latin Modern Sans}
% DejaVu Sans Mono (not Latin Modern Mono) for the monospace face because
% LM Mono lacks Greek glyphs — beta/γ inside lstlisting and \texttt{} were
% silently dropped before this change.
\setmonofont{DejaVu Sans Mono}[Scale=0.85]

\lstset{basicstyle=\ttfamily\scriptsize}

\title{第二讲：计算——Julia 与 VFI}
\subtitle{异质性主体宏观经济学的计算方法}
\author{孙振越 (Jeffrey Sun)}
\institute{多伦多大学}
\date{2026 年 5 月 12 日}

\begin{document}

\maketitle

% =============================================================
\section{准备工作}
% =============================================================

\begin{frame}{贝尔曼方程}
  \[
    \boxed{\;
      V(b, z) \;=\; \max_{c \in [0, b]} \; u(c)
        + \beta \, V(R(b-c) + z,\, z)\quad\forall b,z
    \;}
  \]
  $b$ 是财富。$z$ 是主体的生产率（暂时设为常数）。

  该方程刻画了求解家庭问题所需的函数 $V$。
\end{frame}

\begin{frame}{$V$ 唯一且可计算}
  对于任意有界的 $v_1, v_2$，
  \[
    \norm{\mathcal{T} v_1 - \mathcal{T} v_2}_\infty
      \;\le\; \beta \, \norm{v_1 - v_2}_\infty.
  \]
  非正式地说，巴拿赫不动点定理意味着：
  \begin{itemize}
    \item 存在唯一的 $V^*$ 满足 $V^* = \mathcal{T} V^*$。
    \item 值函数迭代总能找到 $V^*$。
  \end{itemize}
  本讲：用有限网格替代连续状态 $b$，将 $\mathcal{T}$ 写成一个 Julia 函数，迭代至不动点，并绘制 $V$。
\end{frame}

% =============================================================
\section{Julia}
% =============================================================

\begin{frame}[fragile]{数组}
  \begin{lstlisting}
  xs  = [1.0, 2.0, 3.0]      # Vector{Float64}, 长度为 3
  M   = [1 2; 3 4]           # Matrix{Int64}, 2×2
  zs  = zeros(3, 4)          # 3×4 的全 0.0 矩阵
  xs[1]                      # 1.0   --- 从 1 开始索引！
  M[2, 1]                    # 3
  \end{lstlisting}
  \begin{itemize}
    \item \textbf{索引从 1 开始。}
    \item \texttt{xs[end]} 是最后一个元素。\texttt{xs[2:3]} 是一个切片。
    \item 数组按列优先存储：维度 1 是行，维度 2 是列。``向量''是列向量。
  \end{itemize}
\end{frame}

\begin{frame}[fragile]{广播（Broadcasting）}
  \begin{lstlisting}
  xs = [1.0, 2.0, 3.0]

  xs .+ 10          # [11.0, 12.0, 13.0]
  log.(xs)          # [0.0, 0.693, 1.099]
  xs .^ 2           # [1.0, 4.0, 9.0]

  # 对两个形状相同的数组逐元素运算：
  [1,2,3] .* [10,20,30]   # [10, 40, 90]
  \end{lstlisting}
  点号 \texttt{.} 是``逐元素''的标记。

  点号将任意函数或运算逐元素地施加，无需编写循环。
\end{frame}

\begin{frame}[fragile]{函数}
  \begin{lstlisting}
  function utility(c)
      c <= 0 && return -Inf
      return log(c)
  end

  # 简单条件判断的简写形式：
  function utility(c)
      return c <= 0 ? -Inf : log(c)
  end

  # 单行函数同样有效：
  u(c) = c <= 0 ? -Inf : log(c)
  \end{lstlisting}
\end{frame}

% =============================================================
\section{离散化}
% =============================================================

\begin{frame}{离散财富网格}
  \begin{itemize}
    \item $V$ 是定义在 $[0, \infty) \times W$ 上的函数。我们无法精确地表示它。
    \item 选取 $N$ 个财富取值 $b_1 < b_2 < \cdots < b_N$。
    \item 将 $V$ 表示为向量 $[V(b_1), \ldots, V(b_N)]$。
    \item 这使得该函数变为有限维。
  \end{itemize}
  \textbf{注：}值函数迭代的所有概念都\textbf{不依赖}于将 $V$ 表示为向量。$V$ 可以是样条、神经网络，或定义在动态网格上，这些概念仍然成立。
\end{frame}

\begin{frame}[fragile]{线性间距与对数间距}
  \begin{lstlisting}
  loggrid(lo, hi, N) = exp.(range(log(lo), log(hi); length=N))

  # 线性：等间距
  b_grid = collect(range(0.1, 20.0; length=200))

  # 对数：在借贷约束附近更稠密
  b_grid = loggrid(0.1, 20.0, 200)
  \end{lstlisting}
  策略函数在 $b$ 较小时非线性更强，因此我们采用对数间距，在较小的 $b$ 上布置更多网格点。
\end{frame}

\begin{frame}[fragile]{效用}
  \begin{lstlisting}
  u(c) = c <= 0 ? -Inf : log(c)
  \end{lstlisting}
  更一般的形式是带 $\gamma$ 的 CRRA 效用，但现在先保持简单。
\end{frame}

\begin{frame}[fragile]{将贝尔曼算子 $\mathcal{T}$ 写成 Julia 函数}
  \begin{lstlisting}
  function bellman_operator(V, params)
      # 从 `params` 中解包参数
      (; β, R, z, b_grid) = params

      # 计算达到每个下期财富水平 `b_next` 所需的储蓄水平 `b_end`
      b_end = (b_grid' .- z)./R

      # 在 `b_next` 维度（维度 2，即列）上对条件效用 `u.(b_grid .- b_end) .+ β.*V'` 求最大值。
      return maximum(u.(b_grid .- b_end) .+ β.*V'; dims=2) |> vec
  end
  \end{lstlisting}
  \texttt{b\_grid} 是列向量。\texttt{b\_grid'} 是行向量。

  每一行对应一个 $b$ 值；每一列对应一个 $b'$ 值。最后一行在每个 $b$ \textbf{内部}（按行）、\textbf{跨越} $b'$（按列）取 \texttt{max}。
\end{frame}

\begin{frame}{关于示例的说明}
  笔记本中的代码\textbf{并未}经过优化，也不追求速度。它们只是用来阐释概念。

  我在 HouseholdStages.jl 包中提供了快得多的实现。
\end{frame}

% =============================================================
\section{值函数迭代}
% =============================================================

\begin{frame}{值函数迭代}
  \begin{enumerate}
    \item 以 $V_0 \equiv 0$ 起步（任何初始猜测都可以）。
    \vspace{0.4em}
    \item $V_{k+1} \leftarrow \mathcal{T} V_k$。
    \vspace{0.4em}
    \item 当 $\norm{V_{k+1} - V_k}_\infty < \text{tol}$ 时停止。
  \end{enumerate}
\end{frame}

\begin{frame}[fragile]{\texttt{solve\_vfi}}
  \begin{lstlisting}
  function solve_vfi(params; rtol=1e-6, verbosity=0)
      # 给出 V 的初始猜测
      V = zeros(size(params.b_grid))

      # 反复施加贝尔曼算子直至收敛
      Δ = Inf
      while Δ >= rtol
          # 更新 V
          TV = bellman_operator(V, params)

          # 计算误差
          Δ = maximum(abs.(TV .- V))

          # 更新 V
          V = TV
      end

      return V
  end
  \end{lstlisting}
\end{frame}

\begin{frame}[fragile]{基准校准}
  \begin{lstlisting}
  b_grid = loggrid(0.05, 40.0, 400)
  params = (;β=0.96, R=1.04, z=1.0, b_grid)

  V = solve_vfi(params)
  \end{lstlisting}
  \begin{itemize}
    \item \small $\beta R < 1$，否则 $V^*$ 为无穷大。
    \item 网格：从 $0.05$ 到 $40$ 的 400 个对数间距点。
    \item \texttt{params} 是一个 \texttt{NamedTuple}（具名元组）。末尾的 \texttt{b\_grid} 等价于 \texttt{b\_grid=b\_grid}。
  \end{itemize}
\end{frame}

\begin{frame}{值函数的特征}
  $V$ 的形状与 $u$ 的形状相似\footnote{Sheeran (2017)}，
  \begin{itemize}
    \item $V(b)$ 是递增的（财富越多 $\Rightarrow$ 未来消费越多）。
    \item $V(b)$ 是凹的：增加一单位财富，对穷人的价值高于对富人的价值。
    \item 其斜率是\textbf{财富的边际价值}：由包络定理，$V'(b) = u'(c^\star(b))$。
  \end{itemize}
\end{frame}

\begin{frame}[fragile]{策略函数 $c^\star$}
  \[
    c^\star(b) \;=\; \arg\max_{c \in [0, b]} \; u(c) + \beta \, V\!\big(R(b - c) + z\big).
  \]
  \begin{lstlisting}
  function policy_function(V, params)
      (; β, R, z, b_grid) = params

      # 与 bellman_operator 的设置相同
      b_end = (b_grid' .- z) ./ R

      # 在 b'（列）上取 argmax —— 为每一行取出对应的列索引
      idx = argmax(u.(b_grid .- b_end) .+ β .* V'; dims=2)
      return vec(getindex.(idx, 2))::Vector{Int}
  end
  \end{lstlisting}
  与 \texttt{bellman\_operator} 相同，只是用 \texttt{argmax} 取代了 \texttt{maximum}。

  返回一个整数向量：对每个（行）$b$ 给出最优 $b'$ 的网格索引。
\end{frame}

\end{document}
